% concatenated from Separation_Loj_8.tex
\documentclass[reqno, 12pt ,a4letter]{amsart}
\usepackage{amsmath,amsxtra,amssymb,latexsym, amscd,amsthm}
\usepackage[pdftex,pdfpagelabels]{hyperref}
\usepackage{graphicx,color}
%\usepackage[active]{srcltx}
\usepackage[utf8]{inputenc}
\usepackage[mathscr]{eucal}
%\usepackage[mathscr]{euscript}
\usepackage{mathrsfs}
\usepackage[english]{babel}
\usepackage{varioref}
\usepackage{moreverb,graphicx}
\usepackage{enumerate,array}
\usepackage{multirow}
%\usepackage{algorithm}
%\usepackage{algorithm2e}
\usepackage{mathpazo}

%\usepackage{refcheck}
\setlength{\parindent}{12pt}
\setlength{\parskip}{1.0pt}

\setlength{\oddsidemargin}{-.0cm}
\setlength{\evensidemargin}{-.0cm}

\setlength{\textwidth}{6.5in}
\setlength{\textheight}{9in}
\setlength{\headheight}{0in}

\setlength{\topmargin}{-.75cm}
\setlength{\headsep}{1.2cm}
\setlength{\footskip}{1.25cm}

\setlength{\baselineskip}{12pt}
\renewcommand{\baselinestretch}{1.25}

\newtheorem {theorem}{Theorem}[section]
\newtheorem {corollary}[theorem]{Corollary}
\newtheorem {lemma}[theorem]{Lemma}
\newtheorem {claim}{Claim}[section]
\newtheorem {question}{Question}

\theoremstyle{definition}
\newtheorem {definition}[theorem]{Definition}
\newtheorem {remark}[theorem]{Remark}
\newtheorem {example}[theorem]{Example}
\newtheorem {algorithm}{Algorithm}[section]

\newcommand{\ord}{\operatornamewithlimits{ord}}



\def\ees{{\accent"5E e}\kern-.385em\raise.2ex\hbox{\char'23}\kern-.08em}
\def\EES{{\accent"5E E}\kern-.5em\raise.8ex\hbox{\char'23 }}
\def\ow{o\kern-.42em\raise.82ex\hbox{
\vrule width .12em height .0ex depth .075ex \kern-0.16em \char'56}\kern-.07em}
\def\OW{O\kern-.460em\raise1.36ex\hbox{
\vrule width .13em height .0ex depth .075ex \kern-0.16em \char'56}\kern-.07em}

\pagestyle{plain}

\title{On the separation \L ojasiewicz exponents of real analytic sets in the real plane}


\author{PHI-D{U}NG HO{A}NG$^\dag$}
\address{$^\dag$ Department of Mathematics - Faculty of Fundamental Sciences,
\newline \indent  Posts and Telecommunications Institute of Technology,
\newline \indent 
Km10 Nguyen Trai Rd., Ha Dong District, Hanoi, Vietnam}
\email{dunghp@ptit.edu.vn}

\author{HONG-DUC NGUYEN$^\ddag$}
\address{$^\ddag$TIMAS, Thang Long University, Nghiem Xuan Yem, Hanoi, Vietnam}
\email{duc.nh@thanglong.edu.vn}



\thanks{The first and second author's research is funded by Vietnam National Foundation for Science and Technology Development (NAFOSTED) under grant number 101.04-2024.09.}

\keywords{\L ojasiewicz exponents, separation \L ojasiewicz inequality, Newton polygon, Newton--Puiseux roots, effective \L ojasiewicz inequality}

\subjclass{14H20, 14B05, 32B05, 58K05}


%\date{ \today}

\begin{document}
\maketitle

\begin{abstract}
The main aim of the paper is to give a formula for computing the separation \L ojasiewicz exponents for two real analytic set germs via the Newton--Puiseux expansions of their defining functions. Moreover, we present an effective exponent for the case of two real algebraic sets in terms of their degrees.
\end{abstract}

\section{Introduction}
The \L ojasewicz inequalities impressively appeared in the 50s years of last century to solve the question of L. Schwartz for distribution division \cite{Hormander1958,Lojasiewicz1959, Lojasiewicz1965}. They have many relations and applications to other branches of mathematics, such as research on local singularities of analytic functions \cite{Teissier1977}, the proof of Thom's Gradient Conjecture \cite{Kurdyka2000-2}, study of infinitely dimensional version in partial differential equations \cite{Colding2015}, applications in polynomial and tame optimization \cite{HaPham2017,Hoang2016}, \ldots 

Let $A, B$ be real semi-analytic subsets of an open subset $U \subseteq \mathbb{R}^n$ with $0 \in A \cap B$. Then we have the following inequality \cite{Lojasiewicz1965}:
\begin{itemize}
	\item There exist $C, r > 0$ and $\beta \ge 1$ such that 
	\begin{equation}\label{separation-inequality}
		d(x, A) + d(x, B) \geq Cd(x, A \cap B)^\beta, \text{for all}\ \|x\| \leq r,
	\end{equation}
    where $d(x, X)$ is the Euclidean distance of $x \in \mathbb{R}^n$ to the set $X$ ($d(x, X) = 1$ if $X = \emptyset$).
	\item The infimum of such exponents $\beta$ is called {\em separation \L ojasiewicz exponent} of semi-analytic subsets $A$ and $B$ (at the origin), it is denoted by $\mathscr{L}(A,B)$.
\end{itemize} 
Note that the separation \L ojasiewicz exponent $\mathscr{L}(A, B)$ is a rational number \cite[Corollary 2]{Bochnak1975}. The \L ojasiewicz exponent represents and compares the growths of functions on two sides of the inequalities \cite{Lojasiewicz1965}.  

The \L ojasiewicz exponents are used to study some important problems, such as computing some topological invariants \cite{Nguyen2019,Risler1997,Teissier1977}, investigating the singularity at infinity \cite{Ha1990}, the study of the separation of real algebraic sets \cite{Cygan1999,Kurdyka2014} proof of effective Nullstellensatz \cite{Jelonek2005,Ji1992}, etc. There are many works providing computations and estimations of \L ojasiewicz exponents \cite{Gwozdziewicz1999,HD08,HD10,Acunto2005,Kurdyka2014,Nguyen2019,Pham2012}. 
%Recently, the authors in \cite{Dinh2025} give formula of generalized \L ojasiewicz exponent of two functions $f$ and $g$ by using the so called $\rho$-approximations of Newton--Puiseux roots of $f\cdot g = 0$ \cite[Theorem 3.1]{Dinh2025}.

In this paper, we study the separation \L ojasiewicz exponents of the zeros of the real analytic functions in two variables. More precisely, assuming $f$ and $g$ are real analytic functions in two variables, we give a formula to compute the separation \L ojasiewicz exponent $\mathscr{L}(f^{-1}(0),g^{-1}(0))$ in \eqref{separation-inequality} in terms of the approximations of the Newton--Puiseux roots of $f\cdot g$ (see Definition \ref{approximation} and Theorem \ref{separation-theorem}). If $f$ and $g$ are polynomials, then we give an effective formula for the separation \L ojasiewicz exponent. 

%The paper is organized as follows. The theory of Newton polygon relative to an arc and the approximations of Newton--Puiseux roots of a real analytic function germs are presented in Section \ref{Section2}. In Section \ref{Section3}, we give the formula of separation \L ojasiewicz exponent of two real analytic subsets. An effective formula of separation \L ojasiewicz exponent in terms of the degree of two analytic subsets is given in Section \ref{section4}.

\section{Preliminaries} \label{Section2}

\subsection{Newton polygon relative to an arc}\label{subsection-key-lemmas}\quad\\
In this section, we recall the technique of the Newton polygon relative to an arc or sliding-technique due to Kuo-Parusinski \cite{Kuo2000} (see also \cite{HD10}). This technique plays a key role in this article. Let $f: (\mathbb{K}^2,0) \to (\mathbb{K},0)$ ($\mathbb{K}$ is $\mathbb{C}$ or $\mathbb{R}$) be an analytic function germ. Suppose that $f$ is mini-regular in $x$ of order $m$, i.e. in the Taylor expansion of $f$: $$ f(x,y) = f_m(x,y) + f_{m+1}(x,y) + \dots, $$ we have $f_m(1,0) \ne 0$, where $f_k(x,y)$ is the homogeneous component of degree $k$. 

For a {\em Puiseux series}
\begin{eqnarray*}
	x = \phi(y) = c_1y^{n_1/N} + c_2y^{n_2/N}+ \cdots \in \mathbb{K}\{y^{1/N}\}
\end{eqnarray*}
with $N \le n_1 < n_2 < \cdots $ being positive integers, where $c_1 \ne 0$. We define the order of a Puiseux series $x=\phi(y)$:
$$ \mathrm{ord}\phi = \frac{n_1}{N}. $$
The series $\phi$ is said to be real if all coefficients $c_i$ of $\phi$ are real. Let us define 
$$M(X, Y) := f(X + \phi(Y), Y) := \sum c_{ij}X^iY^{j/N}.$$
For each  $c_{ij} \ne 0,$ let us plot a dot at $(i, j/N)$ in $\Bbb R^2$, and call it a {\em Newton dot}. The set of Newton dots is called the {\em Newton diagram}, They generate a convex hull, whose boundary is called the {\em Newton polygon of $f$ relative to $\phi,$}  to be denoted by $\mathbb{P}(f, \phi).$ Note that this is the Newton polygon of $M$ in the usual sense.

%We want to find a Newton-Puiseux root of $f = 0$, i.e. a Puiseux series $x = \gamma(y)$ such that $f(\gamma(y), y) = 0$. The algorithm has the following steps:
\begin{algorithm}[The sliding]\quad
\begin{enumerate}
    \item[\textbf{Input.}]
An analytic function germ $f:(\mathbb{K}^2,0)\to(\mathbb{K},0)$, mini-regular in $x$, and an initial Puiseux series $\phi_0(y)$. Set $k:=0$ and $\phi_0:=\phi$.

%Repeat the following steps for $k=0,1,2,\dots$.
	\item[\textbf{Step 1.}] Define $$M_k(X,Y):=f(X+\phi_k(Y),Y) =\sum c^{(k)}_{ij}X^iY^{j/N}$$ and consider the Newton polygon $\mathbb{P}(f,\phi_k)$ of $M_k$. Let $(0,h_k)$ be the lowest Newton dot on $X=0$ and $E_k$ be the compact edge of $\mathbb{P}(f,\phi_k)$ having $(0,h_k)$ as a vertex. Then $$\ord f(\phi_k(y),y)=h_k.$$	
	\item[\textbf{Step 2.}]
	\begin{itemize}
	    \item If $\mathbb{P}(f,\phi_k)$ has no Newton dot on $X=0$, then $f(\phi_k(y),y)=0$ and the algorithm stops with output $$\phi_\infty(y):=\phi_k(y).$$
        \item If $\mathbb{P}(f,\phi_k)$ has a Newton dot on $X=0$, choose a nonzero root $c_k$ of the edge polynomial
	\[
	\mathcal{E}_{E_k}(z)
	:=\sum_{(i,j/N)\in E_k} c^{(k)}_{ij} z^i
	\] and set
	\[
	\phi_{k+1}(y)
	:=\phi_k(y)+c_k\,y^{\tan\theta_{E_k}},
	\]
	where $\theta_{E_k}$ is the angle of the edge $E_k$.
	\end{itemize}
	\item[\textbf{Step 3.}]
	Increase $k:=k+1$ and return to \textbf{Step 1}.
    \item[\textbf{Output.}] The sequence $\{\phi_k(y)\}_k$ has diagram $\phi_0 \rightarrow \phi_1 \rightarrow \cdots\rightarrow \phi_k \rightarrow \cdots$ which produces a Puiseux series $\phi_\infty$ satisfying
\[
f(\phi_\infty(y),y)=0.
\]
\end{enumerate}
\end{algorithm}
The series $\phi_\infty$ is called a {\it final result of the sliding of $\phi$ along $f$}. %Note that in the iterative process, the lowest dots on $X=0$ of $\mathbb{P}(f,\phi_k)$ will move upward. 
The above algorithm is based on the following lemma.
\begin{lemma} \label{lemma-sliding}
	Suppose that $\phi$ is not a root of $f = 0$. Consider a series 
	$$\psi : x = \phi(y) + c y^{\rho} + o(\rho),$$
	where $c \in \mathbb{K},0<\rho \in \mathbb{Q}$ and $o(\rho)$ denotes a Puiseux series of order greater than $\rho$. Then the following statements hold:
	\begin{itemize}
		\item[(i)] If $\tan \theta_{E_1} < \rho$ or $\tan \theta_{E_1}= \rho$ and $\mathcal{E}_{E_1}(c) \ne 0$ then 
		$\mathbb{P}(f, \phi) = \mathbb{P}(f, \psi),$ and, therefore
		$\mathrm{ord} f(\phi(y), y) = \mathrm{ord} f(\psi(y), y).$
		
		\item[(ii)] If $\tan \theta_{E_1}= \rho$ and $\mathcal{E}_{E_1}(c)  = 0$ then $\mathrm{ord} f(\phi(y), y) < \mathrm{ord} f(\psi(y), y)$.
	\end{itemize}
\end{lemma}
\begin{proof} For a detailed proof, we refer to \cite{HD08}. In fact, the special case where $\psi : x = \phi(y) + c y^{\tan \theta_{E_1}}$ was proved in \cite[Lemma 2.1]{HD08}. The lemma is then deduced by applying the special case (infinitely) many times.
\end{proof}

%Let $f\colon (\mathbb{R}^2,0) \to (\mathbb{R},0)$ be a real analytic function germ, which is mini-regular in $x$ of order $m$. Suppose that $x = \phi(y)$ is a Puiseux series and $x = \beta_j(y)$ are Newton--Puiseux roots of $f(x, y) = 0$ in the ring $\mathbb{R}\{x,y\}$, i.e. real roots of $f=0$. By Lemma \ref{distance}, we put 
We define
$$
\mathrm{ord}d(\phi, V_f) := \begin{cases}\label{order-distance}
	\max\limits_{j} \{\mathrm{ord}(\phi(y)-\beta_j(y))\},\ \text{$\beta_j(y)$ are the real Newton-Puiseux roots of $f$,}\\
	1\ \text{ if $f$ has no real roots,}
\end{cases} 	
$$
where $V_f$ denotes the zero locus $f(x,y)=0$ of $f$.
\begin{lemma}\label{distant-sliding}
Suppose that $\phi$ is not a root of $f = 0$. Then
	\begin{itemize}
		\item[(i)] $\mathrm{ord}d(\phi, V_f)\leq \tan \theta_{E_1}.$		
		\item[(ii)] Let $\psi$ be a series of form
	$$\psi : x = \phi(y) + c y^{\rho} + o(\rho),\ c \in \mathbb{R}.$$
Suppose that $\rho>\mathrm{ord}d(\psi, V_f)$, or $\rho\geq \mathrm{ord}d(\psi, V_f)$ and $c$ is generic. Then
		$$\mathrm{ord}d(\psi, V_f)= \mathrm{ord}d(\phi, V_f).$$
	\end{itemize}
\end{lemma}
\begin{proof}
(i) Assume by contradiction that $\rho:=\mathrm{ord}d(\phi, V_f)> \tan \theta_{E_1}$, then taking a real Newton-Puiseux root $\beta(y)$ such that $\mathrm{ord}d(\phi, V_f)=\mathrm{ord}(\phi(y)-\beta(y))$ we obtain
$$\beta(y)= \phi(y) + c y^{\rho} + o(\rho), c\neq 0.$$
By Lemma \ref{lemma-sliding}, $\mathrm{ord} f(\phi(y), y) = \mathrm{ord} f(\beta(y), y)=+\infty$, which is a contradiction.

(ii) Let $\beta$ be any Newton-Puiseux root of $f$. Then
$$\psi(y)-\beta(y)=\phi(y)-\beta(y)+ c y^{\rho} + o(\rho).$$
Since $\mathrm{ord}\left(\psi(y)-\beta(y)\right)\leq \mathrm{ord}d(\phi, V_f)\leq \rho$, it follows that
$$\mathrm{ord}\left(\psi(y)-\beta(y)\right)=\mathrm{ord}\left(\phi(y)-\beta(y)\right).$$
Therefore
$$\mathrm{ord}d(\psi, V_f)=\mathrm{ord}d(\phi, V_f).$$
\end{proof}
The following lemma is similar to \cite[Proposition 2.3]{HD10}.
\begin{lemma}\label{distance}
	Assume that $\mathrm{ord}\phi(y) \ge 1$. Then
	$$ d( (\phi(y),y), V_f ) \simeq |y|^{\mathrm{ord}d(\phi, V_f)}, \text{ as } y\to 0$$

%Here $A \simeq B$ means that $A/B$ tends to a positive constant.
\end{lemma}
\begin{proof}(cf. \cite[Lemma 2.16]{Kuo1974}). Let us consider the following two cases.
	
	If $f=0$ does not have any real Newton--Puiseux root, then $V_f= (0,0)$. Taking $\epsilon > 0$ small sufficiently such that $|y| < \epsilon$, we have $$d((\phi(y),y),V) = d((\phi(y),y), (0,0)) = \|(\phi(y),y)\| = \sqrt{(\phi(y))^2 + y^2}.$$ Since the assumption $\mathrm{ord}\phi(y) \ge 1$, then 
	\begin{align}\label{Eqn8}
	d((\phi(y),y),V) \simeq |y|^1\ \text{for all}\ |y| < \epsilon.	
	\end{align}
	
	If $f=0$ has some real Newton--Puiseux roots, then there exists $\epsilon > 0$ sufficiently small such that $f(x,y)$ can be represented as follows
	\begin{align*}
		f(x,y) = u(x,y)\prod_{j=1}^{k}(x - \beta_i(y)), 
	\end{align*}
	where $u(x,y)$ has no real roots and $|y| < \epsilon$. Hence
	\begin{align*}
		d( (\phi(y),y), V ) = \inf_j\|(\phi(y),y) - (\beta_j(y), y)\| = \inf_j|\phi(y) - \beta_j(y)|.
	\end{align*}
	It follows that 
	\begin{align*}
		\inf_j|\phi(y) - \beta_j(y)| \simeq |y|^{\max_j\{\mathrm{ord}(\phi(y)-\beta_j(y))\}}.
	\end{align*}
	Therefore, $d((\phi(y),y),V) \simeq |y|^{\mathrm{ord}d(\phi, V)}$  for all $|y| < \epsilon$. Combining this with \eqref{Eqn8}, we deduce the proof of the lemma.
\end{proof}

%%%%%%%%%%%%%%%
\subsection{Approximations of Puiseux series}\label{real-approx}
Let $x = \gamma(y)$ be a Puiseux series in the following form:
\begin{eqnarray*}
	\gamma(y) = a_1 y^{n_1/N} + a_2y^{n_2/N}+ \cdots + a_{s - 1} y^{n_{s - 1}/N} + c_s y^{n_s/N} +  \cdots,
\end{eqnarray*}
where $a_i \in \mathbb{R}$ and $c_s$ is the first non-real coefficient, if there is one. Let us replace $c_s$ by a generic real number $g,$ and call
\begin{equation}\label{real-approximation}
	\gamma^{\mathbb{R}}(y) := a_1 y^{n_1/N} + a_2y^{n_2/N}+ \cdots + a_{s - 1} y^{n_{s - 1}/N} + g y^{n_s/N},
\end{equation}
a {\em real approximation of $\gamma$}. In case $s = +\infty,$ let $\gamma^{\mathbb{R}} = \gamma.$

Let $f\colon (\mathbb{R}^2,0) \to (\mathbb{R},0)$ be a real analytic function germ, which is mini-regular in $x$ of order $m$. Suppose that $\phi$ is not a Newton--Puiseux root of $f = 0$ and $\gamma$ is a final result of the sliding of $\phi$ along $f$. Then
$$\gamma(y)  = \phi(y) + c y^{\tan \theta_{E_1}}+ \text{higher order terms},$$
where $\mathcal{E}_{H_{M}}(c) = 0$ and $M(X,Y) = f(X + \phi(Y),Y)$. If $c \notin \mathbb{R}$, then by the definition of the real approximation \eqref{real-approximation}, the real approximation $\gamma^{\mathbb{R}}$ of the series $\gamma$ also has the form 
\begin{equation}\label{real-approximation2}
	\gamma^{\mathbb{R}}(y) = \phi(y) + g y^{\tan \theta_{E_1}},	
\end{equation}	
where $g \in \mathbb{R}$ is generic. For $f \in \mathbb{K}\{x,y\}$ which is regular in $x$, let $\mathcal{V}_{\mathbb{R}}(f)$ be the set of all real approximations of all the Newton-Puiseux roots of $f$.
\begin{definition}[\cite{Nguyen2019}]
Let $\phi$ be a Puiseux series $\phi(y) = \sum a_iy^{\alpha_i}$. For each positive real number $\rho$, the {\em $\rho$-approximation} of $\phi(y)$ is defined by series 
$$ \sum_{\alpha_i < \rho} a_iy^{\alpha_i} + cy^\rho, $$ where $c$ is a generic real number.
\end{definition}
\begin{remark}
Let $\phi(y)$ be a non-real Puiseux series. Then the real approximation $\phi^{\mathbb{R}}(y)$ of $\phi(y)$ is the $\rho$-approximation of $\phi$, where $\rho$ is the smallest exponent occurring in $\phi$ with non-real coefficient. %It is clear that if $\phi$ is real, then the real approximation of $\phi$ is itself. 
\end{remark}
\begin{definition}\label{approximation}
    Let $\phi_1, \phi_2$ be any two distinct Puiseux series. The series $\phi_{1,2}$ is called the {\em approximation of $\phi_1, \phi_2$} if it is the $\rho$-approximation of $\phi_1$ (also the $\rho$-approximation of $\phi_2$), where $\rho := \text{ord}(\phi_1 - \phi_2)$. Let $f\colon (\mathbb{R}^2,0) \to (\mathbb{R},0)$ be a real analytic function germ, which is mini-regular in $x$. Let $x=\phi_i(y), i=1,\ldots,s$ be the set of the Newton-Puiseux roots of $f$. We denote by $\mathcal{V}_a(f)$ the set of all approximations $\phi_{i,j}$ of $\phi_i$ and $\phi_j$ such that $\phi_{i,j}$ is a real Puiseux series. Note that $\mathcal{V}_{\mathbb{R}}(f) \subset \mathcal{V}_a(f)$ if $f \in \mathbb{R}\{x,y\}$.
\end{definition}

\begin{example}
	Let $\phi_1(y) = y^{\frac{3}{2}} - 2y^{\frac{5}{2}} + 3iy^{\frac{7}{2}} + \dots$ and $\phi_2(y) = y^{\frac{3}{2}} - 2y^{\frac{5}{2}} - 3y^{\frac{7}{2}} + iy^{\frac{9}{2}} + \dots$. Then, the approximation of series $\phi_1$ and $\phi_2$ is $\phi_{1,2}(y) = y^{\frac{3}{2}} - 2y^{\frac{5}{2}} + gy^{\frac{7}{2}}$, where $g \in \mathbb{R}$ and generic.
\end{example}
%%%%%%%%%%%%%%%
\section{Separation \L ojasewicz exponent of two real analytic subsets}\label{Section3}
Let $f, g \colon (\mathbb{R}^2,0) \to (\mathbb{R},0)$ be reduced real analytic functions germs in two variables, which are mini-regular in $x$ of order $m$. Put $V_f = f^{-1}(0), V_g = g^{-1}(0)$. Then, by the \L ojasiewicz inequality \eqref{separation-inequality}, there exist $C, r > 0$ and $\beta \ge 1$ such that
\begin{equation}\label{separation}
	d(x, V_f) + d(x, V_g) \geq Cd(x, V_f \cap V_g)^\beta \text{ for all}\ \|x\| \leq r.
\end{equation} 
%If $x \in V_f$ and $x \not\in V_g$, then there exist $C, r > 0$ and $\beta \ge 1$ such that 

%The infimum of the exponents $\beta$ such that Inequality \eqref{separation} holds, is called the {\em separation \L ojasiewicz exponent} $\mathscr{L}(V_f, V_g)$.

Let $\varphi\colon [0,\varepsilon)\to \Bbb R^2$ be a semi-algebraic analytic curve such that $\varphi(0)=0$. Then there is a rational number $\ell(\varphi)$ such that
$$d(\varphi(t), V_f)+d(\varphi(t), V_g)\sim \left(d(\varphi(t), V_f \cap V_g)\right)^{\ell(\varphi)}\text{ as } t\to 0.$$
Let $\mathscr{L}(V_f, V_g)$ denote the separation \L ojasiewicz exponent of $V_f$ and $V_g$. Using the Curve Selection Lemma (see \cite{Milnor1968}), we can easily prove the following equality
\begin{eqnarray}\label{Eqn10}
	\mathscr{L}(V_f, V_g) &=& \sup_\phi \ell(\phi),
\end{eqnarray}
where the supremum is taken over all semi-algebraic analytic curves passing through the origin, which do not contain in $V_f \cap V_g$.

Suppose that the curve $\varphi$ is parametrized by a Puiseux series $x= \phi(y)$. We can see that
\begin{equation}\label{lphi-separation}
	\ell(\varphi)=\ell(\phi) := \frac{ \min\{\mathrm{ord}d(\phi, V_f), \mathrm{ord}d(\phi, V_g)\} }{\mathrm{ord}d(\phi,V_{\gcd(f,g)})}.	
\end{equation} 
We define the number
\begin{align}\label{L+_formula}
   \mathscr{L}_+(V_f, V_g) := \max\{\ell(\gamma)| \gamma \in \mathcal{V}_a(f \cdot g)\},
\end{align}
where the set $\mathcal{V}_a(f \cdot g)$ is defined in Definition \ref{approximation}. Let us denote by $\overline{f}, \overline{g}$ the real analytic function germs defined by $\overline{f}(x, y) := f(x, -y)$ and $\overline{g}(x, y) := g(x, -y)$, and set $\mathscr{L}_{-}(V_f, V_g) := \mathscr{L}_{+}(V_{\overline{f}}, V_{\overline{g}})$. Now, we establish the formula for the separation \L ojasiewicz exponent $\mathscr{L}(V_f, V_g)$.
\begin{theorem}\label{separation-theorem}
	Let $f, g \colon (\mathbb{R}^2,0) \to (\mathbb{R},0)$ be non-zero reduced real analytic function germs, which are mini-regular in $x$ of order $m$. Then
	$$\mathscr{L}(V_f, V_g) = \max\{\mathscr{L}_{+}(V_f, V_g), \mathscr{L}_{-}(V_f, V_g)\},$$ where $\mathscr{L}_{-}(V_f, V_g) := \mathscr{L}_{+}(V_{\overline{f}}, V_{\overline{g}})$ and $\mathscr{L}_{+}(V_{\overline{f}}, V_{\overline{g}})$ is defined by \eqref{L+_formula}.
\end{theorem}
First, we need some lemmas. Put $h := \gcd(f,g)$, where $\gcd(f,g)$ is the greatest common division of $f$ and $g$. Let us consider the Newton polygons relative to any arc $x=\phi(y)$, they are $\mathbb{P}(f, \phi)$, $\mathbb{P}(g,\phi)$ and $\mathbb{P}(h,\phi)$ with the highest Newton edges $F_1$, $G_1$ and $H_1$, respectively. Assume that $\mathcal{E}_{F_1}$, $\mathcal{E}_{G_1}$ and $\mathcal{E}_{H_1}$ are their associated polynomials. Let $\theta_{F_1}$, $\theta_{G_1}$ and $\theta_{H_1}$ be the Newton angles of $F_1$, $G_1$ and $H_1$, respectively. It is easy to see that 
	\begin{align}\label{compare-3-tan}
		\min\{\tan\theta_{F_1}, \tan\theta_{G_1}\} \geq \tan\theta_{H_1}.
	\end{align}
	\begin{lemma}\label{claim1} 
		Suppose that $\ell(\phi) > \mathscr{L}_+(V_f, V_g)$. Then 
    \begin{itemize}
        \item[(i)] $\tan\theta_{F_1} = \tan\theta_{G_1}.$
        \item[(ii)] The polynomial $\mathcal{E}_{F_1}\mathcal{E}_{G_1}$ has only one root, which is a real value.
    \end{itemize}
		\end{lemma}
	\begin{proof}
    %It is easy to see that $\phi$ is the form 	$$ (x = \phi(t), y = t)\ \text{or}\ (x = \phi(t), y = -t)\ \text{where}\ \phi(t) \in \mathbb{R}\{t^{1/N}\}. $$ 	Without loss of generality, assume that $\phi$ can be parametrized by $(x = \phi(t), y = t)$ and $\tan\theta_{F_1} > \tan\theta_{G_1}$. 
    By contradiction we assume that $\tan\theta_{F_1} \neq \tan\theta_{G_1}.$ Without loss of generality suppose that $\tan\theta_{F_1} > \tan\theta_{G_1}.$ Let $\phi_{1,\infty}$ and $\phi_{2,\infty}$ be final results of sliding $\phi$ along $f$ and $g$, respectively. Then we have 
	\begin{align}
		\phi_{1,\infty}(y) &= \phi(y) + \sum_{i \ge 1}a_iy^{\alpha_i}, a_1\neq 0\\
		\phi_{2,\infty}(y) &= \phi(y) + \sum_{i \ge 1}b_iy^{\beta_i}, b_1\neq 0,
	\end{align}
	where $\tan\theta_{F_1}=\alpha_1 <\alpha_2 < \ldots$ and $\tan\theta_{G_1} =\beta_1 < \beta_2 <
    \ldots $.
    
    Let $\phi_{1,2}$ be the approximation of $\phi_{1,\infty}$ and $\phi_{2,\infty}$. Then $\phi_{1,2}\in \mathcal{V}_a(fg)$ and
	$$ \phi_{1,2}(y) = \phi(y) + cy^{\tan\theta_{G_1}} + o(\tan\theta_{G_1}), $$ where $c \in \mathbb{R}$ is generic. It follows from Lemma \ref{distant-sliding} that
    	\begin{align}\label{order-with-Vg}
    		\mathrm{ord}d(\phi_{1,2}, V_g) = \mathrm{ord}d(\phi, V_g) \text{ and } \mathrm{ord}d(\phi_{1,2}, V_h) = \mathrm{ord}d(\phi, V_h).
    	\end{align}
    	 
We will show that $\ell(\phi_{1,2})=\ell(\phi)$ by considering the following two cases.

\begin{itemize}
	\item Case 1: 	$\mathrm{ord}d(\phi, V_f) \leq \tan\theta_{G_1}$. Then $\mathrm{ord}\ d(\phi_{1,2}, V_f) = \mathrm{ord}\ d(\phi, V_f)$ due to Lemma \ref{distant-sliding}. Hence $\ell(\phi_{1,2})=\ell(\phi)$.
    
	\item Case 2: $\mathrm{ord}d(\phi, V_f) > \tan\theta_{G_1}$. Take a real Newton-Puiseux root $\gamma$ of $f$ satisfying
	\begin{align*}
		\mathrm{ord}(\phi - \gamma) = \mathrm{ord}\ d(\phi, V_f).
	\end{align*}
    Since
    $$ \phi_{1,2}(y)- \gamma(y) =\phi(y)-\gamma(y) + cy^{\tan\theta_{G_1}} + o(\tan\theta_{G_1}), $$
     it follows that $\mathrm{ord}d(\phi_{1,2} - \gamma) = \tan\theta_{G_1}$. Therefore,
	\begin{align}\label{order-with-Vf-tan}
		\mathrm{ord}d(\phi_{1,2}, V_f) \ge \tan\theta_{G_1}\ge \mathrm{ord}d(\phi_{1,2}, V_g).
	\end{align}
	Hence,
	  		\begin{align*}
			\ell(\phi_{1,2}) &= \frac{\min\{\mathrm{ord}d(\phi_{1,2},V_f), \mathrm{ord}d(\phi_{1,2},V_g)\}}{\mathrm{ord}d(\phi_{1,2},V_h)} \\
			&= \frac{ \mathrm{ord}d(\phi,V_g) }{\mathrm{ord}d(\phi,V_h)}  =\ell(\phi).
		\end{align*} 
		
\end{itemize}
The equality $\ell(\phi_{1,2})=\ell(\phi)$ contradicts the assumption that $\ell(\phi) > \mathscr{L}_+(V_f, V_g)$.  This proves (i).

(ii) By (i) one has $\rho=\tan \theta_{F_1} = \tan\theta_{G_1}$. Then $\mathcal{E}_{F_1}\mathcal{E}_{G_1}$ is the polynomial associated with the highest Newton edge of $\mathbb{P}(fg, \phi)$ with the angle $\rho$. Suppose by contradiction that $\mathcal{E}_{F_1}\mathcal{E}_{G_1}$ has two distinct roots $c_1,c_2$. Let $\varphi_i=\phi+c_i y^{\rho }$ for each $i=1,2$, and let $\varphi_{i,\infty}$ be a final result of sliding $\varphi_i$ along $fg$. Let $\varphi_{1,2}$ be the approximation of $\varphi_{1,\infty}$ and $\varphi_{2,\infty}$. Then $\varphi_{1,2}\in \mathcal{V}_a(fg)$ and 
	$$ \varphi_{1,2}(y) = \phi(y) + cy^{\rho} + o(\rho), $$ where $c \in \mathbb{R}$ is a generic number. Note that 
    $$\tan \theta_{H_1}\leq \tan \theta_{F_1}=\tan \theta_{G_1}=\rho.$$
    It follows from Lemma \ref{distant-sliding} that
    $$\mathrm{ord}d(\varphi_{1,2},V_f)=\mathrm{ord}d(\phi,V_f),\ \mathrm{ord}d(\varphi_{1,2},V_g)=\mathrm{ord}d(\phi,V_g) $$
	and
    $$d(\varphi_{1,2},V_h)=\mathrm{ord}d(\phi,V_h).$$
    Hence $\ell(\phi_{1,2})=\ell(\phi)$, which is a contradiction. The lemma is proved.    
\end{proof}

\begin{proof}[Proof of Theorem \ref{separation-theorem}]
	By the formula \eqref{Eqn10}, it is obvious that $$ \mathscr{L}(V_f, V_g) \geq \max\{\mathscr{L}_{+}(V_f, V_g), \mathscr{L}_{-}(V_f, V_g)\}. $$
	
	Suppose for contradiction that $$ \mathscr{L}(V_f, V_g) > \max\{\mathscr{L}_{+}(V_f, V_g), \mathscr{L}_{-}(V_f, V_g)\}. $$
	Then there is a real analytic arc $\phi$ passing through the origin and not lying the $x$-axis such that 
	\begin{align}\label{contradicted-assumption}
	 \ell(\phi) > \max\{\mathscr{L}_{+}(V_f, V_g), \mathscr{L}_{-}(V_f, V_g)\}	
	\end{align}
	and the parameter forms of $\phi$ is 
	$$ (x = \phi(t), y = t)\ \text{or}\ (x = \phi(t), y = -t),\ \text{where}\ \phi(t) \in \mathbb{R}\{t^{1/N}\}. $$
    Without loss of generality, assume that $\phi$ can be parametrized by $(x = \phi(t), y = t)$. 
    By Lemma \ref{claim1}, then $\theta_{F_1} = \theta_{G_1}$ and the polynomial $\mathcal{E}_{F_1}\mathcal{E}_{G_1}$ has a unique root $c_0 \in \mathbb{R}$. Define 
$$ \tilde\phi(y) := \phi(y) + c_0y^{\tan\theta_{F_1}}. $$
\begin{claim}\label{claim4} We have
	$$\ell(\tilde\phi) \ge \ell(\phi).$$
\end{claim}
\begin{proof}
We first see that $$\min\{\mathrm{ord}d(\phi, V_f),\  \mathrm{ord}d(\phi, V_g)\} > \mathrm{ord}d(\phi, V_f \cap V_g)=\mathrm{ord}d(\phi, V_h),$$
because $\ell(\phi)>1$. Then $$\tan\theta_{F_1}\geq \mathrm{ord}d(\phi, V_f)> \mathrm{ord}d(\phi, V_h).$$
Applying Lemma \ref{distant-sliding} and its proof, one has
$$\mathrm{ord}d(\tilde\phi, V_h)=\mathrm{ord}d(\phi, V_h),\ \mathrm{ord}d(\tilde\phi, V_f)\geq \mathrm{ord}d(\phi, V_f) \text{ and } \mathrm{ord}d(\tilde\phi, V_g)\geq\mathrm{ord}d(\phi, V_g).$$
Hence $\ell(\tilde\phi) \ge \ell(\phi).$
\end{proof}

Consider Newton polygons $\mathbb{P}(f,\tilde\phi)$ and $\mathbb{P}(g,\tilde\phi)$. The polygon $\mathbb{P}(f,\tilde\phi)$ has the highest Newton edge $\widetilde{F}_1$ with angle $\theta_{\widetilde{F}_1}$ and the associated polynomial $\mathcal{E}_{\widetilde{F}_1}$; the polygon $\mathbb{P}(g,\tilde\phi)$ has the highest Newton edge $\widetilde{G}_1$ with angle $\theta_{\widetilde{G}_1}$ and the associated polynomial $\mathcal{E}_{\widetilde{G}_1}$. Combining Claim \ref{claim4} and Lemma \ref{claim1} we obtain 
%%%%%%%

\begin{claim}\label{claim3}
	If $\tilde\phi$ is not a Newton--Puiseux root of $h=0$, then
	\begin{enumerate}
		\item $ \tan\theta_{\widetilde{F}_1} = \tan\theta_{\widetilde{G}_1}$;
		\item The polynomial $\mathcal{E}_{\widetilde{F}_1}\mathcal{E}_{\widetilde{G}_1}$ has only real root.
	\end{enumerate}
\end{claim}
%%%%%%%%%%%

%%%%%%%%%
Let $\phi_0:=\phi$ and $$\phi_1(y)=\tilde\phi(y) = \phi_0(y) + c_0y^{\tan\theta_{F_1}}.$$ 
Then $\phi_1$ is the unique sliding of $\phi_0$ along $f$ (as well as $g$ and $h$). Applying this process, we obtain a sequence 
\begin{align}
    \phi \to \gamma_1 \to \gamma_2 \to \dots \to \gamma_n \to \dots \to \phi_\infty,
\end{align}
where $\phi_{k+1}$ be the unique sliding of $\phi_k$ along $h$ for each $k\geq 1$. The series $\phi_\infty$ is a real Newton-Puiseux root of $h$. It yields the following identities
$$\mathrm{ord}d(\phi, V_f)=\mathrm{ord}d(\phi, V_g)=\mathrm{ord}d(\phi, V_h)=\mathrm{ord}(\phi-\phi_\infty)=\tan \theta_{F_1}.$$
This implies $\ell(\phi)=1$, which is a contradiction. So, the theorem is proved.
\end{proof}
%%%%%%%%%%%%%%%

\section{Effective estimation for separation \L ojasiewicz exponent}\label{section4}
Suppose that $f$ and $g$ are polynomials of two real variables, then there is the following effective version of the separation \L ojasiewicz exponent.
\begin{theorem}
    Let $f$ and $g$ be polynomials of two real variables such that their degrees are not greater than $d$ and $f(0)=g(0)=0$. Then there exist positive numbers $C$ and $r$ such that 
    \begin{equation}\label{effective-separation}
	d(x, V_f) + d(x, V_g) \geq Cd(x, V_f \cap V_g)^{\frac{(2d-1)^2+1}{2}},\ \text{for all}\ \|x\| \leq r.
\end{equation}
\end{theorem}
\begin{proof}(cf. \cite{Kurdyka2014})
Put $F := f^2 + g^2$, then $\deg F \le 2d$, where $d = \max\{\deg f, \deg g\}$. Moreover, $V_F = V_f \cap V_g$. 

If $V_f \subset V_g$, then $\mathscr{L}(V_f, V_g) = 1$. Assume that $V_f \not\subset V_g$. Applying the classical \L ojasiewicz inequality to $F$, there exists $\alpha, c > 0$ such that
$$ |F(x)| \ge  c d(x, V_F)^\alpha, \forall x \in U \cap (V_f \setminus V_g),$$
where $U \subset \mathbb{R}^2$ is a neighbourhood of $(0,0)$. 
Let $\widetilde{\mathcal{L}}(F)$ be the classical \L ojasiewicz exponent of the polynomial $F$, then, by \cite[Theorem 4.4]{Nguyen2019}, $\widetilde{\mathcal{L}}(F) \le (2d-1)^2+1$. Hence, the effective optimal version of classical \L ojasiewicz inequality is 
$$ |F(x)| \ge  c d(x, V_F)^{(2d-1)^2+1}, \forall x \in U \cap (V_f \setminus V_g), $$
Note that $x \in U \cap (V_f \setminus V_g)$, then
\begin{align}\label{F and g}
|F(x)| = |g(x)|^2.    
\end{align}
Therefore,
\begin{align}
   |g(x)| \ge  c d(x, V_f \cap V_g)^{\frac{(2d-1)^2+1}{2}}, \forall x \in U \cap (V_f \setminus V_g),
\end{align}
It is easy to see that there exists $C, C' > 0$ such that
\begin{align}\label{sim of d and g}
   Cd(x, V_g) \ge C'|g(x)|. 
\end{align}
Combining \eqref{F and g} with \eqref{sim of d and g}, we obtain 
\begin{align*}
    d(x, V_g) \ge Cd(x, V_f \cap V_g)^{\frac{(2d-1)^2+1}{2}}, \forall x \in U \cap (V_f \setminus V_g).
\end{align*}
This implies inequality \eqref{effective-separation}.
\end{proof}
%%%%%%%%%
\begin{remark}
The exponent in inequality \eqref{effective-separation} is sharper than the exponent in \cite[Corollary 8]{Kurdyka2014} in the case of two variables, where $0 \in \mathbb{R}^2$ is not necessarily isolated point of $V_f \cap V_g$. 
\end{remark}
%%%%%%%%%%%%%%%
%\subsection*{Acknowledgment.} The first and second author's research is funded by Vietnam National Foundation for Science and Technology Development (NAFOSTED) under grant number 101.04-2024.09.


\bibliographystyle{amsalpha}
\begin{thebibliography}{AA}

\bibitem{Bochnak1975} 
J. Bochnak and J. -J. Risler, {\em Sur les exposants de {{\L}}ojasiewicz}, Comment. Math. Helv., {\bf 87}(4) (1975), 493--507.

\bibitem{Brieskorn1986} 
E. Brieskorn and H. Kn\H{o}rrer, {\em Plane algebraic curves}, Birkh\H{a}user Verlag, Basel, 1986.

\bibitem{Colding2015}
T. B. Colding; W. P. Minicozzi II, {\em Uniqueness of blowups and Łojasiewicz inequalities}, Ann. of Math. (2) 182 (2015), no. 1, 221–285.

\bibitem{Cygan1999}
E. Cygan, T. Krasiński, and P. Tworzewski, {\em Separation of algebraic sets and the Łojasiewicz exponent of polynomial mappings}, Invent. Math. 136 (1999), no. 1, 75–87.

\bibitem{Acunto2005} 
D. D'Acunto and K. Kurdyka, {\em Explicit bounds for the {{\L}}ojasiewicz exponent in the gradient inequality for polynomials}, Ann. Polon. Math., {\bf 87} (2005), 51--61. 

\bibitem{Dinh2025}
S. T. Dinh, F. Guo, H. D. Nguyen, T. S. Pham, {\em Computation of the Łojasiewicz exponents of real bivariate analytic functions}, Manuscripta Math. 176, 1 (2025), 21 pp.

\bibitem{Gwozdziewicz1999} 
J. Gwo\'zdziewicz, {\em The \L ojasiewicz exponent of an analytic function at an isolated zero,} Comment. Math. Helv. Vol. {\bf 74}(3) (1999), 364--375.

\bibitem{Ha1990}
H. V. Ha, {\em Nombres de \L ojasiewicz et singulariti\'{e}s \`{a} l'infini des polyn\^{o}mes de deux variables complexes}, C. R. Acad. Sci. Paris, S\'{e}ries I Math. 311 (1990), 429-432.

\bibitem{HD08} 
H. V. H\` a and H. D. Nguyen, {\em On the \L ojasiewicz exponent near the fibre of polynomial mappings,} Ann. Polon. Math. {\bf 94} (2008), 43--52.

\bibitem{HD10} 
H. V. H\` a and H. D. Nguyen, {\em \L ojasiewicz inequality at infinity for polynomials in two real variables,} Math. Z. {\bf 266} (2010), 243--264.

\bibitem{HaPham2017}
H. V. H\` a and T. S. Pham, {\em Genericity in polynomial optimization}, Ser. Optim. Appl. 3. World Scientific, Singapore, 2017.

\bibitem{Hoang2016}
P. D. Hoang, {\em \L ojasiewicz-type inequalities and global error bounds for nonsmooth definable functions in o-minimal structures}, Bull. Aust. Math. Soc., 93 (2016), no. 1, 99–112.

\bibitem{Hormander1958}
L. H\"ormander, {\em On the division of distributions by polynomials,} Ark. Mat. {3} N. 53 (1958), 555--568.

\bibitem{Jelonek2005}
Z. Jelonek, {\em On the effective Nullstellensatz}, Invent. Math. 162 (1) (2005) 1–17.

\bibitem{Ji1992}
S. Ji, J. Kollar and B. Shiffman, {\em A global Lojasiewicz inequality for algebraic varieties}, Trans. Amer. Math. Soc. 329(2) (1992), 813–818.

\bibitem{Kuo1974} 
T. C. Kuo, {\em Computation of {{\L}}ojasiewicz exponent of $f(x, y)$}, Comment. Math. Helv., {\bf 49} (1974), 201--213.

\bibitem{Kuo2000} 
T. C. Kuo and A. Parusi\'nski, {\em Newton polygon relative to an arc,} Real and Complex Singularities (S\~ao Carlos, 1998), Chapman \& Hall Res. Notes Math., {\bf 412} (2000), 76--93.

\bibitem{Kurdyka2014}
 K. Kurdyka and S. Spodzieja, {\em Separation of real algebraic sets and the {{\L}}ojasiewicz exponent}, Proc. Amer. Math. Soc., {\bf 142}(9) (2014), 3089--3012.

\bibitem{Kurdyka2000-2}
K. Kurdyka, T. Mostowski, A. Parusiński, {\em Proof of the gradient conjecture of R. Thom}, Ann. of Math., {\bf 152}(3) (2000), 763--792.

\bibitem{Lojasiewicz1959}  
S. \L ojasiewicz, {\em Sur le probl\`eme de la division}, Studia Math.{\bf 18} (1959), 87--136.

\bibitem{Lojasiewicz1965} 
S. {{\L}}ojasiewicz, {\em Ensembles semi-analytiques}, Inst. Hautes Etudes Sci., Bures-sur-Yvette, France, 1964.

\bibitem{Milnor1968} 
J. W. Milnor, {\em Singular points of complex hypersurfaces,} Annals of Mathematics Studies vol. {\bf 61}, Princeton Univ. Press, USA, 1968.

\bibitem{Nguyen2019}
H. D. Nguyen,  T. S. Pham and P. D. Hoang, {\em Topological invariants of plane curve singularities: Polar quotients and Lojasiewicz gradient exponents}, Internat. J. Mathematics, 30 (14) (2019), 1950073, 19 pp.

\bibitem{Pham2012}
T. S. Pham, {\em An explicit bound for the {{\L}}ojasiewicz exponent of real polynomials}, Kodai Math. J., {\bf 35}(2) (2012), 311--319.

\bibitem{Risler1997}
J.-J. Risler and D. Trotman, {\em Bi-Lipschitz invariance of the multiplicity,} Bull. London. Math. Soc., {\bf 29}(2) (1997), 200--204.

\bibitem{Teissier1977}
B. Teissier, {\em Vari\'et\'es polaires. I. Invariants polaires des singularit\'es d'hypersurfaces,} Invent. Math. {\bf 40}(3) (1977), 267--292.
\end{thebibliography}
\end{document}
%%%%%%%%%%%%%
\section{On the \L ojasiewicz exponent of the generalized \L ojasiewicz inequality}\label{Section4}
In the paper \cite{Dinh2025}, the authors gave two formulas of the generalized \L ojasiewicz exponent $\mathscr{L}_g(f)$. The first formula for $\mathscr{L}_g(f)$ is represented by the quotients of the multiplicities of roots of $f=0, g=0$ and the so-called $\rho$-approximations of roots of $f \cdot g = 0$ \cite[Theorem 3.1]{Dinh2025}. The second formula for $\mathscr{L}_g(f)$ is represented by the quotients of the multiplicities of roots and the real approximations of roots of $f=0$ \cite[Theorem 3.2]{Dinh2025}. From the idea of Ha \cite[Theorem 2.8]{Ha2018}, that is $\mathcal{L}(f)$ attained along the real polar curve $f_x=0$. Hence, Question \ref{question1} appeared, we recall it. 

{\em Can the \L ojasiewicz exponent $\mathscr{L}_g(f)$ be attained along the real polar curve $f_x=0$?}

The following theorem give the answer for above question.
\begin{theorem}\label{counterexample}
	Let $f, g \colon \mathbb{R}^2 \to \mathbb{R}$ be real analytic functions with $f(0,0) = g(0,0) = 0$, which is mini-regular in $x$ of order $m$. Then the \L ojasiewicz exponent $\mathscr{L}_g(f)$ can not be attained along the real polar curve $f_x = 0$.
\end{theorem}
\begin{proof}
	Consider the following polynomials $$ f(x) = x^2 + xy + y^2,\ g(x) = (2x+y)^2 + y^{10}. $$ 
	We have $\overline{f}(x,y) = x^2 - xy + y^2$ and $g(x) = (2x - y)^2 + y^{10}$.
	
	It is easy to see that two Newton--Puiseux roots of $f(x,y)=0$ is 
	\begin{align*}
		x = \left(\frac{-1 + i\sqrt{3}}{2}\right)y\quad \text{and}\quad x = \left(\frac{-1 - i\sqrt{3}}{2}\right)y. 
	\end{align*}
	and two Newton--Puiseux roots of $\overline{f}(x,y) = 0$ is
	\begin{align*}
		x = \left(\frac{1 + i\sqrt{3}}{2}\right)y\quad \text{and}\quad x = \left(\frac{1 - i\sqrt{3}}{2}\right)y. 
	\end{align*}
	Hence the real approximations of two Newton--Puiseux roots of $f$ are $x = \varphi_1(y) = c_1y$ and $x = \varphi_2(y) = c_2y$, where $c_1, c_2$ are generic real numbers. Hence 
	\begin{align*}
		\mathrm{ord}f(\varphi_1)= \mathrm{ord}f(\varphi_2) = 2, \mathrm{ord}g(\varphi_1) = \mathrm{ord}g(\varphi_1) = 2.
	\end{align*}
	Similarly, we can compute for $\overline{f}$ and $\overline{g}$. Since \cite[Theorem 3.2]{Dinh2025}, then we have $\mathscr{L}_g^+(f) = \mathscr{L}_g^-(f) = 1.$ Therefore,
	$$ \mathscr{L}_g(f) = 1. $$
	
	Let us consider the polar curve $\frac{\partial f}{\partial x} = 0$, that is $2x + y = 0$. Its root is $x = \phi(y) = -\frac{y}{2}$. Then $\mathrm{ord}f(\phi) = 2$ and $\mathrm{ord}g(\phi) = 10$. Similarly, the real polar curve $\frac{\partial \overline{f}}{\partial x} = 2x - y = 0$ has a root $x = \overline{\phi}(y) = \frac{y}{2}$. Then $\mathrm{ord}\overline{f}(\overline{\phi}) = 2$ and $\mathrm{ord}\overline{g}(\overline{\phi}) = 10$. 
	
	Therefore $$l(\phi) = l(\overline{\phi}) = \frac{2}{10} < \mathscr{L}_g(f) = 1.$$ So we can conclude that the \L ojasiewicz exponent $\mathscr{L}_g(f)$ can not be attain along the polar curve $f_x = 0$.
\end{proof}

	\begin{align}
	\min\{\mathrm{ord}d(\phi, V_f), \mathrm{ord}(\phi, V_g)\} \leq \min\{\tan\theta_{F_1},\tan\theta_{G_1}\} \leq \min\{\mathrm{ord}(\phi^\mathbb{R}, V_f), \mathrm{ord}(\phi^\mathbb{R},V_g\}.	
\end{align}
